\documentclass[a4paper]{article}
\usepackage[utf8]{inputenc}
\usepackage[slovene]{babel}
\usepackage{graphicx}
\usepackage{hyperref}
\usepackage[nottoc]{tocbibind}
\usepackage{minted}
\usepackage{listings}
\usepackage{caption}
\usepackage{subcaption}
\usepackage{amsmath}
\usepackage{ dsfont }
\usepackage{siunitx}
\usepackage{multimedia}
\usepackage[table,xcdraw]{xcolor}
\setlength\parindent{0pt}

\definecolor{codegreen}{rgb}{0,0.6,0}
\definecolor{codegray}{rgb}{0.5,0.5,0.5}
\definecolor{codepurple}{rgb}{0.58,0,0.82}
\definecolor{backcolour}{rgb}{0.95,0.95,0.92}
\newcommand{\ddd}{\mathrm{d}}
\newcommand\myworries[1]{\textcolor{red}{#1}}
\newcommand{\Dd}[3][{}]{\frac{\ddd^{#1} #2}{\ddd #3^{#1}}}

\lstdefinestyle{mystyle}{
    backgroundcolor=\color{backcolour},   
    commentstyle=\color{codegreen},
    keywordstyle=\color{magenta},
    numberstyle=\tiny\color{codegray},
    stringstyle=\color{codepurple},
    basicstyle=\ttfamily\footnotesize,
    breakatwhitespace=false,         
    breaklines=true,                 
    captionpos=b,                    
    keepspaces=true,                 
    numbers=left,                    
    numbersep=5pt,                  
    showspaces=false,                
    showstringspaces=false,
    showtabs=false,                  
    tabsize=2
}

\lstset{style=mystyle}

\begin{document}
\begin{titlepage}
    \begin{center}
        \includegraphics[]{logo.png}
        \vspace*{3cm}
        
        \Huge
        \textbf{Reševanje PDE z metodo Galerkina}
        
        \vspace{0.5cm}
        \large
        11. naloga pri Matematično-fizikalnem praktikumu

        \vspace{4.5cm}
        
        \textbf{Avtor:} Marko Urbanč (28191096)\ \\
        \textbf{Predavatelj:} prof. dr. Borut Paul Kerševan\ \\
        
        \vspace{2.8cm}
        
        \large
        8.9.2023
    \end{center}
\end{titlepage}
\tableofcontents
\newpage
\section{Uvod}
Če poznamo lastne funkcije diferencialnega operatorja za določeno geometrijo, se reševanje parcialnih
diferencialnih enačb včasih lahko prevede na razvoj po lastnih funkcijah. V tem primeru se lahko
diferencialni operator zapiše kot matrika in enačbo potem rešujemo kot sistem linearnih enačb. Tega lahko
računamo kot vemo in znamo. Zdaj smo to počeli že parkrat.\\

V našem primeru bo fizikalna inspiracija Navier-Stokesova enačba, ki je pravzaprav drugi Newtonov zakon za 
tekočine. Vendar pa je ta enačba zelo zapletena in je še vedno odprt problem, ali sploh obstajajo rešitve
v splošnem. Zato se bomo omejili na preprostejši primer, kjer privzamemo, da imamo enakomeren laminaren tok 
nestisljive tekočine v dolgi ravni cevi pod vplivom stalnega tlačnega gradienta $p'$. V tem primeru se 
Navier-Stokesova enačba poenostavi na Poissonovo enačbo

\begin{equation}
    \nabla^2 \vec{v} = -\frac{p'}{\eta}\>,
    \label{eq:poisson}
\end{equation}

kjer je $\vec{v}$ hitrost tekočine in $\eta$ njena viskoznost. Enačbo rešujemo v notranjosti preseka cevi, medtem ko
je ob stenah hitrost enaka nič. Za pretok velja Poiseuillov zakon

\begin{equation}
    \Phi = \int_S{v\>\ddd S}  = C \frac{p' S^2}{8\pi\eta}\>,
    \label{eq:poiseuille}
\end{equation}

kjer je $S$ presek cevi in $C$ konstanta, ki je odvisna od oblike preseka. Konstanta znaša $C = 1$ za krožni presek.
V našem primeru pa bomo določili konstanto $C$ pa polkrožni presek. Uvedemo nove spremenljivke $\xi = r/R$ in 
$u=v\eta/(p'R^2)$ in nato se problem glasi

\begin{equation}
    \Delta u(\xi,\>\phi) = -1, \qquad u(1,\>\phi) = u(\xi,\>0) = u(\xi,\>\pi) = 0\>,
    \label{eq:poisson2}
\end{equation}
\begin{equation}
    C = 8\pi \iint{\frac{u(\xi,\>\phi)\xi\>\ddd\xi\> \ddd\phi}{(\pi/2)^2}}\>.
\end{equation}

Da se izognemo računanju lastnih funkcij (v temu primeru Besselovih) in njihovih ničel, lahko zapišemo
rešitev v obliki razvoja po neki poskuni bazi. V našem primeru bomo vzeli bazo

\begin{equation}
    \psi_{nm}(\xi,\>\phi) = \xi^{2m+1}(1-\xi)^n \sin((2m+1)\phi)\>.
\end{equation}

Z njo lahko zapišemo aproksimacijo rešitve kot linearno kombinacijo 

\begin{equation}
    \tilde{u}(\xi,\>\phi) = \sum_{n=0}^{\infty}\sum_{m=0}^{\infty}{a_{nm}\psi_{nm}(\xi,\>\phi)}\>.
\end{equation}

Za te funkcije niti ni nujno da so prava baza v smislu, da bi bile ortogonalne druga na drugo. Potrebno je
da zadoščajo robnim pogojem tako, da jim bo avtomatično zadoščala tudi linearna kombinacija. Približna funkcija 
$\tilde{u}$ seveda ne zadosti Poissonovi enačbi za res, ampak ji preostane majhna napaka $\varepsilon$

\begin{equation}
    \Delta \tilde{u}(\xi,\>\phi) = -1 + \varepsilon(\xi,\>\phi)\>.
\end{equation}

Pri metodi Galerkina zahtevamo, da je napaka ortogonalna na vse poskusne funkcije $\psi_{nm}$, torej

\begin{equation}
    \langle \psi_{nm},\>\varepsilon \rangle = 0 \qquad \forall n,m\>.
\end{equation}

V splošnem bi lahko zahtevali še ortogonalnost napake na nek drug sistem utežnih funkcij. Metoda Galerkina je
poseben primer takih metod (angl. \textit{Methods of Weighted Residuals}), kjer je utežna funkcija kar poskusna
funkcija sama. Omenjena izbira vodi do sistema enačb za koeficiente $a_{nm}$

\begin{equation}
    {A_{nm,n'm'} a_{n'm'}} = b_{nm}\>.
\end{equation}

Koeficiente $b_{nm}$ dobimo iz skalarnega produkta

\begin{equation}
    b_{nm} = \langle -1,\>\psi_{nm} \rangle\>,
\end{equation}

ki se zaradi ortogonalnosti poskusnih funkcij poenostavi v

\begin{equation}
    b_{nm} = - \frac{2}{2m + 1} \mathrm{B}(2m+3, n+1)\>,
\end{equation}

kjer je $\mathrm{B}$ Eulerjeva beta funkcija. Matrika $A$ pa je definirana kot 

\begin{equation}
    A_{nm,n'm'} = \langle \nabla^2\psi_{nm},\>\psi_{n'm'} \rangle\>,
\end{equation}

kar se po upoštevanju ortogonalnosti poskusnih funkcij poenostavi v

\begin{equation}
    A_{nm,n'm'} = -\frac{\pi}{2} \frac{nn'(3+4m)}{2+4m+n+n'} \mathrm{B}(n+n'-1, 3+4m)\>\delta_{mm'}\>.
\end{equation}

Končno, se naša enačba za koeficient za pretok $C$ glasi

\begin{equation}
    C = -\frac{32}{\pi}\sum_{mn,m'n'}{b_{nm}A^{-1}_{nm,n'm'}b_{m'n'}}\>.
    \label{eq:C}
\end{equation}

\section{Naloga}
Naloga od nas zahteva, da rešimo Poissonovo enačbo (\ref{eq:poisson2}) in izračunamo koeficient za 
pretok $C$ v primeru polkrožnega preseka cevi. Naj si tudi pogledamo kako je odvisna natančnost
rešitve od števila členov v indeksih $m$ in $n$.\\
\section{Opis reševanja}

\section{Rezultati}


\section{Komentarji in izboljšave}


\newpage
\bibliographystyle{unsrt}
\bibliography{sources}
\end{document}
